\item Un contenedor ligero, cúbico de volumen $a^3$ inicialmente está lleno de un líquido de densidad de masa $\rho$ como se muestra en la figura P15.51a. El cubo es soportado inicialmente por una cuerda ligera para formar un péndulo simple de longitud $L_i$, medido desde el centro de la masa del recipiente lleno, donde $L_i \gg a$. El líquido puede fluir desde el fondo del recipiente a razón constante ($dM/dt$). En cualquier tiempo $t$, el nivel del líquido en el recipiente es $h$ y la longitud del péndulo es $L$ (medida con respecto del centro de masa instantáneo) como se muestra en la figura P15.51b.
\begin{enumerate}
    \item Encuentre el periodo del péndulo en función del tiempo.
    \item ¿Cuál es el periodo del péndulo después de que el contenedor se vacía por completo?
\end{enumerate}